\section{Closures at a fixed corner}\label{cl:section}

Combining cuts can improve a bound even when every individual cut is
suboptimal. The relevant question is therefore whether the intersection of
all cuts from a family recovers the corner-hull dominant. Throughout this
section the vertex and the projected rays remain fixed. Cuts obtained after
reoptimization belong to a different closure.

For a nonempty family $\mathcal F$ of closed convex $S$-free sets containing
$\sbar$ in their interior, put
\[
 V=\{a(C):C\in\mathcal F\}\subseteq\R_+^N,\qquad
 \mathcal C_{\mathcal F}=\{\lambda\ge0:a^T\lambda\ge1\quad(a\in V)\}.
\]
Write $z_{1,\mathcal F}(\rc)=\sup_{C\in\mathcal F}z_C(\rc)$ and
$z_{\mathrm{cl},\mathcal F}(\rc)=\inf_{\lambda\in\mathcal C_{\mathcal F}}
\rc^T\lambda$. Validity gives
\begin{equation}\label{cl:basic-bounds}
 \mathcal D\subseteq\mathcal C_{\mathcal F},\qquad
 z_{1,\mathcal F}\le z_{\mathrm{cl},\mathcal F}\le\zK.
\end{equation}
For the one-cut expression we use positive costs, so no convention for
$0\cdot\infty$ is needed.

\begin{lemma}[Coefficient description]\label{cl:blocking}
Let $\widehat V=\cl(\conv V+\R_+^N)$ and $U=\cl(V+\R_+^N)$.
For $z>0$ and $\rc\ge0$,
\[
 z_{\mathrm{cl},\mathcal F}(\rc)\ge z
 \quad\Longleftrightarrow\quad \rc/z\in\widehat V.
\]
For $\rc>0$, $z_{1,\mathcal F}(\rc)\ge z$ if and only if
$\rc/z\in U$. Consequently, the closure bound equals the best single-cut
bound for every positive cost vector if and only if $U$ is convex.
\end{lemma}

\begin{proof}
Every vector of $\widehat V$ has scalar product at least one with every
point of $\mathcal C_{\mathcal F}$. If $\rc/z\notin\widehat V$, strict
separation gives $\mu^T\rc/z<\beta\le\mu^Ta$ for all
$a\in\widehat V$. Upward closure forces $\mu\ge0$, and therefore
$\beta>0$. The point $\mu/\beta$ belongs to $\mathcal C_{\mathcal F}$ and
has cost less than $z$.

For a single cut, $z_C(\rc)\ge z$ means $a(C)\le\rc/z$.
A supremum of at least $z$ therefore gives, for each $\epsilon\in(0,1)$,
a vector $a\in V$ with $a\le\rc/((1-\epsilon)z)$.
Conversely, if $\rc/z=\lim(a_n+d_n)$ with $d_n\ge0$, positivity of
$\rc$ implies $a_n\le\rc/((1-\epsilon)z)$ for all sufficiently large $n$.
These observations prove the characterization by $U$. If $U$ is convex,
then $U=\widehat V$. If the two bounds agree for all positive costs, every
strictly positive point of $\widehat V$ belongs to $U$; adding
$\epsilon\mathbf1$ and taking limits proves $\widehat V=U$.
\end{proof}

The next theorem identifies the cuts that can contribute to an exact bound.
It permits limits of sets and does not require attainment of the best
single-cut value.

\begin{theorem}[Exactness from tight limit cuts]\label{cl:tight}
Let $S$ be closed, $\rc>0$, and $z=\zK(\rc)\in(0,\infty)$.
Then $z_{\mathrm{cl},\mathcal F}(\rc)=z$ if and only if there exist
$m\le N$, vectors $a^{(1)},\ldots,a^{(m)}\in\cl V$, and
weights $\theta_i>0$ with $\sum_i\theta_i=1$ such that
\begin{equation}\label{cl:dominated-combination}
 \sum_i\theta_i a^{(i)}\le\rc/z.
\end{equation}
Whenever this holds, every selected vector is tight at every corner
minimizer $\lambda^*$, and the slack in
\cref{cl:dominated-combination} is orthogonal to every such minimizer.
\end{theorem}

\begin{proof}
The reverse implication follows from \cref{cl:blocking} and validity.
For the forward implication, \cref{cl:blocking} supplies
\[
 b_n=\sum_{i=1}^{N+1}\theta_{n,i}a_{n,i}+d_n\longrightarrow\rc/z,
 \quad a_{n,i}\in V,\quad d_n\ge0,\quad
 \theta_{n,i}\ge0,\quad\sum_i\theta_{n,i}=1,
\]
where finite-dimensional Carath\'eodory gives $N+1$ slots.
Take a subsequence on which every weight converges, say to $\theta_i$.
If $\theta_i>0$, nonnegativity gives
$0\le a_{n,i}\le b_n/\theta_{n,i}$ componentwise, so that slot is bounded.
Pass to a further subsequence on which each such vector converges to
$a^{(i)}\in\cl V$. Slots with zero limiting weight can have unbounded
vectors. Dropping their nonnegative contributions preserves the required
inequality. The retained limiting weights still sum to one.

Positive costs and closedness ensure the existence of a minimizer in $X$.
Limits of valid cut vectors remain valid, so
\[
 1=\rc^T\lambda^*/z
 \ge\sum_i\theta_i(a^{(i)})^T\lambda^*
 \ge\sum_i\theta_i=1.
\]
Both inequalities are equalities. Every selected vector is tight, and
the slack is orthogonal to $\lambda^*$. This argument applies to every
minimizer. The selected vectors lie in the affine hyperplane
$a^T\lambda^*=1$, of dimension at most $N-1$, because $\lambda^*\ne0$.
Carath\'eodory within that hyperplane reduces their number to at most $N$
without changing their convex combination.
\end{proof}

\begin{proposition}[A smooth active face]\label{cl:smooth}
Suppose $S=\{s:\widetilde g(s)\le0\}$, where $\widetilde g$ is $C^1$
near the projected corner. Let $\rc>0$, $z=\zK(\rc)\in(0,\infty)$,
and let $\lambda^*$ be a minimizer with support $J$. Put
$g(\lambda)=\widetilde g(\sbar+P\lambda)$ and assume
$\nabla_Jg(\lambda^*)\ne0$.
Every valid linear inequality $a^T\lambda\ge1$ tight at $\lambda^*$
satisfies $a_J=\rc_J/z$.

Let $T=\{a\in\cl V:a_J=\rc_J/z\}$. Then
\begin{equation}\label{cl:face-characterization}
 z_{\mathrm{cl},\mathcal F}(\rc)=z
 \quad\Longleftrightarrow\quad
 \rc_{J^c}/z\in\conv\{a_{J^c}:a\in T\}+\R_+^{J^c}.
\end{equation}
At most $|J^c|+1$ tight limit vectors are needed on the right.
\end{proposition}

\begin{proof}
The minimizer satisfies $g(\lambda^*)=0$: if the value were negative,
a small move toward zero would preserve feasibility and lower the cost.
On the relative interior of the face supported on $J$, the implicit
function theorem makes $g=0$ a smooth hypersurface near $\lambda^*$.
For every tangent vector $d$ there is a curve in that hypersurface with
derivative $d$. The valid tight inequality has a local minimum on that
curve, so $a_J^Td=0$. Thus $a_J$ is a multiple of
$\nabla_Jg(\lambda^*)$. The same applies to the valid tight inequality
$\rc^T\lambda/z\ge1$. Evaluation at $\lambda^*$ fixes both multiples:
their scalar products with $\lambda^*_J$ equal one. Hence
$a_J=\rc_J/z$.

Necessity of \cref{cl:face-characterization} follows from
\cref{cl:tight}. For sufficiency, any convex combination on the right,
extended by its fixed coordinates on $J$, is dominated by $\rc/z$;
\cref{cl:tight} applies. Carath\'eodory in $\R^{J^c}$ gives the stated
number of vectors.
\end{proof}

\begin{corollary}[At most one unused ray]\label{cl:one-unused}
Under the hypotheses of \cref{cl:smooth}, if $|J^c|\le1$, then
\[
 z_{\mathrm{cl},\mathcal F}(\rc)=\zK(\rc)
 \quad\Longleftrightarrow\quad z_{1,\mathcal F}(\rc)=\zK(\rc).
\]
\end{corollary}

\begin{proof}
All selected tight limit cuts agree with $\rc/z$ on $J$.
If no coordinate remains, any one of them equals $\rc/z$.
If one coordinate remains, a convex average below its target has at least
one constituent below that target. Thus some $a\in\cl V$ satisfies
$a\le\rc/z$, and \cref{cl:blocking} gives best-single-cut exactness.
The converse follows from \cref{cl:basic-bounds}.
\end{proof}

For $q=w-xy$ and three projected rays spanning $\R^3$, a support-two
minimizer satisfies the regularity condition in \cref{cl:smooth}.
Indeed, if its two independent directions were tangent to $q=0$, they
would span the tangent plane. Projection of this plane onto $(x,y)$ is
an isomorphism, and the quadratic term $-d_xd_y$ is indefinite there.
At a stationary feasible contact, a negative quadratic direction and its
opposite are both locally feasible. Minimality would force the positive
objective to vanish on an open set of directions, a contradiction.
Thus \cref{cl:one-unused} transfers every strict single-cut obstruction
in a support-two direction to the closure. It does not assert equality
of the two suboptimal bounds.

\begin{proposition}[A dimension bound on the gain]\label{cl:gain}
For $\rc>0$ with $z_{\mathrm{cl},\mathcal F}(\rc)\in(0,\infty)$,
\[
 z_{1,\mathcal F}(\rc)\ge z_{\mathrm{cl},\mathcal F}(\rc)/N.
\]
The factor $N$ is attained by nonnegative coefficient families.
\end{proposition}

\begin{proof}
Put $z=z_{\mathrm{cl},\mathcal F}(\rc)$.
By \cref{cl:blocking}, $\rc/z$ belongs to the boundary of
$\widehat V$: membership of $\rc/z'$ for $z'>z$ would contradict the
definition of $z$. A supporting normal $\mu\ne0$ can be chosen
nonnegative, since $\widehat V$ is upward closed.
The bounded-aggregate argument in the proof of \cref{cl:tight} gives
$b=\sum_i\theta_i a^{(i)}\le\rc/z$, with $a^{(i)}\in\cl V$.
Support and nonnegativity imply
\[
 \mu^T\rc/z\le\mu^Tb\le\mu^T\rc/z.
\]
Each selected vector therefore lies in the supporting hyperplane.
Its dimension is at most $N-1$, so at most $N$ vectors suffice.
Some weight is at least $1/N$, and then
$a^{(i)}\le N\rc/z$. Apply \cref{cl:blocking}.
For $V=\{e_1,\ldots,e_N\}$ and $\rc=\mathbf1$, the one-cut bound is one
and the closure bound is $N$.
\end{proof}

\subsection{Bilinear families and exact counterexamples}\label{cl:bilinear}

We use the orbit family $\famA$ and the completion family $\famB$ of
\cref{fd:orbit,fd:completion}. A completion may contain the vertex in its
interior even when its uncompleted orbit set does not.
The transformed point rule gives a smaller family, denoted $\famBP$.
Writing $\Xi=F^TM(\sbar)$, its orbit parameters are exactly
$\Xi=\Sigma\succ0$ symmetric; $\famBP$ contains their upward completions.
The transformations here are all determinant-preserving homogeneous
automorphisms, including transposition. They are not restricted to
permutations. The parameter derivation appears in
\cref{cl:point-parameters}.
For comparison, let $\mathcal P$ denote the corresponding uncompleted
point-rule family. Containment of sets and inclusion of families give
\begin{equation}\label{cl:family-inclusions}
 \mathcal C_{\famB}\subseteq\mathcal C_{\famA}
 \subseteq\mathcal C_{\mathcal P},\qquad
 \mathcal C_{\famB}\subseteq\mathcal C_{\famBP}
 \subseteq\mathcal C_{\mathcal P}.
\end{equation}

\begin{theorem}[Strict closure loss]\label{cl:counterexample}
Let
\[
 \sbar=(-9/2,0,3/2),\quad
 v_1=(-1,-6,18),\quad v_2=(-5,6,-18),\quad
 v_3=(0,5/2,5/2),\quad p_j=v_j-\sbar.
\]
Then $\zK(\mathbf1)=1$, with unique minimizer $(1/2,1/2,0)$,
and the closures of $\famA,\famB,\famBP$, and $\mathcal P$ all strictly
contain $\mathcal D$. In particular,
\begin{align*}
 &(67/250,41/100,151/500)\in\mathcal C_{\famB},\qquad
 (3/11,41/99,10/33)\in\mathcal C_{\famB},\\
 &(51/350,0,17/150)\in\mathcal C_{\famBP},\qquad
 (1/5,0,1/6)\in\mathcal C_{\famBP}.
\end{align*}
In the direction $\mathbf1$ the certified brackets are
\[
 0.97538\le\zA\le z_{\mathrm{cl},\famA}
 \le z_{\mathrm{cl},\famB}\le49/50,
 \qquad
 2539/10000\le\zBP\le z_{\mathrm{cl},\famBP}\le136/525.
\]
The same lower bound $0.97538$ holds for $\zB$.
\end{theorem}

\begin{proof}
The corner geometry and the strict inequality $\zB(\mathbf1)<1$ are
proved in \cref{fd:counterexample}. The support-two minimizer is regular:
$P_{\{1,2\}}^T\nabla q=(-3/2,-3/2)$ at its image.
By \cref{cl:one-unused}, exactness of any of the four closures would imply
best-single-cut exactness, contradicting \cref{fd:counterexample} and
\cref{cl:family-inclusions}.
The displayed points and numerical brackets have exact finite
certificates described in \cref{cl:finite-certificates}.
Every point of $\mathcal D$ has coordinate sum at least one, whereas the
four displayed points have sums $49/50,98/99,136/525$, and $11/30$.
\end{proof}

The larger orbit example in \cref{fd:counterexample} has
$\sbar=(-1/2,1/2,1)$, endpoints
$(1/2,-4,-1/2)$, $(-10,3,24)$, and $(9/2,1/2,15/2)$, and a
support-two minimizer with unit costs. Its strict $\famA$ single-cut
loss therefore also proves strict $\famA$ closure loss. Its numerical
$\famB$ loss is not an exact certificate for that family.

\begin{proposition}[Strict improvement with support one]\label{cl:support-one}
At the corner
\[
 \sbar=(-2,3,2),\qquad v_1=(0,0,0),\quad
 v_2=(6,-2,1/4),\quad v_3=(1,-5/2,1/2),\quad p_j=v_j-\sbar,
\]
unit costs give a unique corner minimizer $e_1$ and
\[
 0.9838\le\zA\le0.9839
 <0.9953611\le z_{\mathrm{cl},\famA}
 \le999/1000<1=\zK.
\]
\end{proposition}

\begin{proof}
The corner value and the single-cut bracket are proved in
\cref{fd:support-one}. Three rational orbit sets and a rational LP dual
prove the closure lower bound; an exact finite certificate proves
$(24/25,1/100,29/1000)\in\mathcal C_{\famA}$.
All data and certificate implications appear in
\cref{cl:support-one-certificate,cl:finite-certificates}.
\end{proof}

Thus the closure can be strictly better than every single cut and still
fail to be exact. The lower bound also holds for $\famB$ by
\cref{cl:family-inclusions}. The upper bound $999/1000$ has only been
certified for $\famA$; the analogous $\famB$ upper bound remains a
numerical observation.

\subsection{Finite and unbounded approximation factors}\label{cl:factors}

For a fixed corner with $X\ne\varnothing$, define
\[
 \rho_{\mathcal F}
 =\sup_{\rc>0}\frac{\zK(\rc)}{z_{\mathrm{cl},\mathcal F}(\rc)}.
\]
Equivalently, it is the least $\rho\ge1$ such that
$\mathcal C_{\mathcal F}\subseteq(1/\rho)\mathcal D$, with value infinity
if there is no such $\rho$. Strict positive-cost separation proves this
equivalence, as well as
$\mathcal C_{\mathcal F}=\mathcal D$ if and only if all positive-cost
bounds are exact. The gain factor in \cref{cl:gain} compares closure and
single-cut bounds; $\rho_{\mathcal F}$ compares the closure with the
corner bound.

\begin{theorem}[Unbounded fixed-corner factors]\label{cl:unbounded}
Let $\epsilon>0$ in every cost vector below.
At the corner of \cref{cl:counterexample}, every $\famBP$ cut has
$a_1\ge7/2$, and this bound is attained. Hence
$(2/7,0,0)\in\mathcal C_{\famBP}$ and, for $\rc_\epsilon=(\epsilon,1,1)$,
\[
 z_{\mathrm{cl},\famBP}(\rc_\epsilon)\le2\epsilon/7,
 \qquad \zK(\rc_\epsilon)\ge367/759.
\]
In particular $\rho_{\famBP}=\rho_{\mathcal P}=\infty$.

Now take the vertex $(1,-1,1)$ and projected rays
$e_x,-e_y,e_w,-e_w$. They are the projections of the independent rays
$e_x,-e_y,e_w+e_z,-e_w+e_z$ in four variables. At this corner,
\[
 \mathcal D=\{\lambda\ge0:\lambda_4\ge2\},\qquad
 \zK(\rc)=2\omega_4.
\]
For $\rc_\epsilon=(\epsilon,\epsilon,\epsilon,1)$,
\begin{align*}
 &z_{\mathrm{cl},\famA}(\rc_\epsilon)\le60\epsilon,
 &&\rho_{\famA}=\infty,\\
 &z_{\mathrm{cl},\famBP}(\rc_\epsilon)\le4(\sqrt2+1)\epsilon,
 &&\rho_{\famBP}=\infty,\\
 &\mathcal C_{\famB}=\mathcal D,
 &&\rho_{\famB}=1.
\end{align*}
Moreover, every $\famBP$ parameter $\Sigma\succ0$ satisfies
\[
 a_1+a_2\ge(\sqrt2-1)/2,
\]
with equality exactly when $\Sigma$ is a positive multiple of $I$.
Consequently,
$(t,t,0,0)\in\mathcal C_{\famBP}$ exactly when $t\ge2(\sqrt2+1)$.
\end{theorem}

\begin{proof}
Exact witnesses and full proofs are in
\cref{cl:bp-factor,cl:wcorner-proof}.
At the four-ray corner, expansion gives
$q=2+\lambda_1+\lambda_2+\lambda_1\lambda_2+\lambda_3-\lambda_4$.
The point $2e_4\in X$ and the inequality $\lambda_4\ge2$ give the stated
dominant. An integer matrix certificate puts $(20,20,20,0)$ in
$\mathcal C_{\famA}$. The $\famBP$ coefficient inequality and its
attainment put $(2(\sqrt2+1),2(\sqrt2+1),0,0)$ in
$\mathcal C_{\famBP}$. The completion with $F^T=J^T$,
$J=\left(\begin{smallmatrix}0&1\\-1&0\end{smallmatrix}\right)$, is
\[
 \{x\ge0,\ y\le0,\ w\ge1-2\sqrt{-xy}\},
\]
and gives the exact cut $\lambda_4/2\ge1$.
All assertions of unboundedness follow by letting $\epsilon$ decrease
to zero in the positive-cost directions displayed above.
\end{proof}

At the four-ray corner \cref{dp:depth-bound} also gives
$(2-\sqrt2)\epsilon\le\zA(\rc_\epsilon)$ for
$0<\epsilon\le\sqrt2$. This auxiliary range is necessary: the bound
$\zA\le\zK=2$ precludes that linear lower estimate for all
$\epsilon>0$. The corner is removed immediately by bound propagation
if a solver knows $x\ge1$, $y\le-1$, and $w\le1$; it establishes a
geometric distinction between the families.

\begin{proposition}[A finite-factor class]\label{cl:finite-factor}
If every projected ray first meets $S$ at a finite step $t_j$, then any
member $C\in\mathcal F$ gives
\[
 \rho_{\mathcal F}\le\max_j\frac{t_j}{\alpha_j(C)}<\infty.
\]
\end{proposition}

\begin{proof}
The first-hit points belong to $X$, so
$\zK(\rc)\le\min_j\omega_jt_j$. Since $C$ is $S$-free,
$0<\alpha_j(C)\le t_j<\infty$; and
\[
 z_{\mathrm{cl},\mathcal F}(\rc)
 \ge\min_j\omega_j\alpha_j(C)
 \ge\min_j\frac{\alpha_j(C)}{t_j}\min_j\omega_jt_j.
\]
\end{proof}

\begin{corollary}[No uniform factor across corners]\label{cl:no-uniform}
There is no positive uniform lower bound for
$z_{\mathrm{cl},\famA}/\zK$ or $z_{\mathrm{cl},\famB}/\zK$, even with
three rays, positive unit costs, and a fixed vertex $(0,0,1)$.
The same conclusion holds for $\famBP$.
\end{corollary}

\begin{proof}
The three-ray construction in \cref{dp:vanishing} has
$\zB/\zK\to0$. By \cref{cl:gain},
$z_{\mathrm{cl},\famB}/\zK\le3\zB/\zK\to0$.
The rescaling $(x,y,w)\mapsto(x/\sqrt\epsilon,y/\sqrt\epsilon,w/\epsilon)$
maps its vertex to $(0,0,1)$ and preserves the corner values and family
ratios. The inclusion relations give the conclusions for $\famA$ and
$\famBP$. If the quantitative single-cut estimate
$\zB/\zK\le137\sqrt\epsilon/z_0$, with
$z_0=(1+\sqrt{1+4\epsilon})/2$ and $0<\epsilon<1$, is used, the closure
estimate is $z_{\mathrm{cl},\famB}/\zK\le411\sqrt\epsilon/z_0$.
The quantitative estimate is \cref{dp:certified-sharpness};
the qualitative conclusion needs only \cref{dp:vanishing}.
\end{proof}

At the three-ray corner of \cref{cl:counterexample}, a further exact
certificate gives $\rho_{\famA}\ge\rho_{\famB}>1.1166$;
see \cref{cl:local-factor}. Whether either fixed-corner factor is finite
there remains open. Numerical estimates near $1.128$ do not establish
finiteness.

\subsection{The unrestricted closure and later corners}\label{cl:unrestricted}

By \cref{fd:dominant}, all $S$-free sets together recover
$\mathcal D$ at one corner. Thus the preceding fixed-corner gaps result
from restricting the set family. This equality of sets gives exact
unrestricted closure bounds for nonnegative costs as well as positive
costs; it does not extend positive-cost single-cut sufficiency or
attainment to zero costs.

The rank-one closure of a polyhedron uses its original bases. It need not
equal the hull of its nonlinear feasible points, even when the
unrestricted family is used at every basis. The disk example in
\cref{fd:rank-one-example} proves this and closes its gap at a newly
created basis in the second round. The distinction is geometric: each
original basis contributes the dominant of its own corner, whereas the
polyhedron has additional constraints outside that corner.

The bilinear records exhibit the same distinction. Reoptimization with
one orbit cut per round reached the original corner bound after three
cuts at \cref{cl:counterexample} and after two at
\cref{cl:support-one}. These are numerical trajectories, not finite
convergence theorems. The complete retained survey and its one-sided
bounds appear in \cref{cl:numerical-records}.

\citet{ABP2018} study constant-factor approximation by lattice-free
closures at a fixed apex. Under their closure and target-representation
hypotheses, approximation of a target cut by the entire family is
equivalent, up to constants, to approximation by an individual member.
Our smooth-face criterion concerns exactness in one objective direction
and uses the support of a corner minimizer. Its hypotheses do not give
a direct application of their theorem to the quadratic families.
The coefficient description and the dimension bound above are elementary
finite-dimensional aggregation statements; the bilinear obstructions and
certified support-one improvement identify what happens for the
restricted quadratic families.
